\section{Discussion}
\label{sec:discussion}

% Outline:
%
% 6.1 Implications for supply chain resilience practice
%   - The methodology (Section 3) + tooling (Section 4) together close the
%     research-to-practice gap that the SCR literature repeatedly flags.
%   - Audit-ready optimisation as a precondition for operational adoption.
%
% 6.2 Transferability beyond humanitarian allocation
%   - The same FIS+EA+production-tooling pattern applies to other SCR
%     decisions: supplier-portfolio diversification under disruption,
%     last-mile rerouting, multi-tier inventory under uncertainty.
%   - The presidio-hardened-vol-assign architecture (pluggable indices,
%     objectives, solvers) is designed for this transfer.
%
% 6.3 Limitations
%   - FIS rule bases require expert elicitation; subjective.
%   - EAs are stochastic; require multiple runs for confidence.
%   - Test cases are illustrative; real disaster response demands integration
%     with live data feeds (out of scope for v0.2.0).
%
% 6.4 Threats to validity
%   - Construct: do FLPP/TFL/CABL capture the resilience constructs we claim?
%   - Internal: parameter sensitivity addressed in Section 5.3.
%   - External: generalisability beyond the two test cases.

TODO.
